% Industry CV — one page, for robotics / perception / ML engineering roles.
% Build: xelatex CV_Industry.tex
\documentclass[10pt]{article}
\usepackage{cvstyle}

\begin{document}

% ============================================================ header
{\fontsize{22}{26}\selectfont\bfseries Hong-Son Nguyen}\\[3pt]
{\fontsize{10.5}{13}\selectfont\color{muted} Robotics Engineer — autonomous navigation, perception, reinforcement learning}\\[8pt]
{\fontsize{9.4}{14}\selectfont
\ilink{\faEnvelope}{hongsonnguyen.haui@gmail.com}{mailto:hongsonnguyen.haui@gmail.com} \quad
\itext{\faPhone*}{+84 866 401 750} \quad
\itext{\faMapMarker*}{Tainan, Taiwan}\\[2pt]
\ilink{\faGlobe}{sontypo.github.io}{https://sontypo.github.io} \quad
\ilink{\faGithub}{sontypo}{https://github.com/sontypo} \quad
\ilink{\faLinkedin}{hong-son-nguyen}{https://www.linkedin.com/in/hong-son-nguyen-751tps/} \quad
\ilink{\aiGoogleScholar}{Google Scholar}{https://scholar.google.com/citations?user=FENrtWsAAAAJ}}

% ============================================================ summary
\section{Profile}
{\fontsize{9.7}{12.4}\selectfont
Robotics researcher who ships complete autonomy stacks, not just models. Three
years building navigation systems that run on real hardware — perception,
prediction, planning, learning and control — for mobile robots operating around
people. Seven publications including IEEE ICRA 2026.}

% ============================================================ skills first
\section{Technical skills}

\skillrow{Languages}{Python, \cpp}
\skillrow{ML}{PyTorch, deep reinforcement learning (PPO, RSL-RL), graph neural networks, selective state space models (Mamba)}
\skillrow{Robotics}{ROS~2, Nav2, NVIDIA Isaac Sim / IsaacLab, Gazebo, motion planning, Social Force Model}
\skillrow{Perception}{3D LiDAR (Velodyne VLP-16), stereo vision (ZED-X), pedestrian detection and tracking}
\skillrow{Hardware}{Agilex Scout Mini, Jetson AGX Orin, edge deployment}

% ============================================================ experience
\section{Experience}

\entry{Feb 2024 -- present}{Research Fellow — Robotics}{Networked Robotic Systems Laboratory, National Cheng Kung University, Taiwan}
\begin{achievements}
  \item Own an end-to-end human-aware navigation stack running on an Agilex
        Scout Mini: LiDAR and stereo perception, pedestrian tracking, trajectory
        prediction, planning and control, deployed on Jetson AGX Orin.
  \item Designed Social Graph Mamba, a trajectory forecasting model that replaces
        quadratic attention with selective state space models over dynamic
        interaction graphs, giving linear scaling in the number of tracked people
        while matching benchmark accuracy on ETH/UCY and SDD.
  \item Built large-scale RL training environments in Isaac Sim / IsaacLab with
        configurable crowds, obstacles and curricula, and trained navigation
        policies with PPO via RSL-RL.
  \item Designed and ran human-subject experiments with real pedestrians,
        turning comfort surveys and their statistical analysis into an empirical
        model used by the planner. Four papers in two years, one at IEEE ICRA 2026.
\end{achievements}

\entry{2022 -- Jan 2024}{Research Assistant, Team Leader}{Hanoi Intelligent Robotics Laboratory, Hanoi University of Industry, Vietnam}
\begin{achievements}
  \item Led a student research team on reinforcement learning for autonomous
        mobile robots, coordinating the work that produced the group's first
        published navigation papers.
  \item Implemented and benchmarked Q-Learning, SARSA and a deep RL / Pure-Pursuit
        hybrid controller in Gazebo, with validation on physical robots.
\end{achievements}

% ============================================================ projects
\section{Selected projects}

\entry{2026}{Social Graph Mamba}{Attention-free multi-agent trajectory prediction \ \plink{arXiv}{https://arxiv.org/abs/2608.21411}}
\begin{detail}
Selective state space models over dynamically built interaction graphs, with a
community-aware module using differentiable MinCut to recover pedestrian groups.
Linear sequence complexity instead of the quadratic cost of attention, evaluated
on ETH/UCY and SDD and deployed on a physical robot.
\end{detail}

\entry{2025}{RL Social Navigation}{Large-scale training environment in NVIDIA IsaacLab}
\begin{detail}
Configurable crowds, obstacles, curricula and goal distributions for learning
socially aware navigation policies, trained with PPO through RSL-RL and
transferred to the Scout Mini platform.
\end{detail}

\entry{2023}{Autonomous Hexapod Robot}{Undergraduate thesis — graded 94, selected as best thesis}
\begin{detail}
Deep reinforcement learning for locomotion and outdoor autonomy on an 18-DOF
hexapod, from simulation training through to hardware deployment.
\end{detail}

% ============================================================ education
\section{Education}

\entry{2024 -- present}{M.S. in Robotics and Control Systems}{National Cheng Kung University, Taiwan — GPA 3.52/4.0}
\entry{2019 -- 2023}{B.S. in Mechatronic Engineering}{Hanoi University of Industry, Vietnam — GPA 3.05/4.0}

% ============================================================ extras
\section{Publications, awards and languages}

\skillrow{Publications}{Seven papers on robot navigation and pedestrian prediction, including IEEE ICRA 2026 and two Springer volumes — \href{https://scholar.google.com/citations?user=FENrtWsAAAAJ}{Google Scholar}.}
\skillrow{Awards}{Distinguished International Student Scholarship, National Cheng Kung University (2024--2026); first prize, Students Scientific Research Competition, Hanoi University of Industry (2022); departmental scholarship (2023); ASME student member since 2022.}
\skillrow{Languages}{Vietnamese (native); English — IELTS 6.0}

\end{document}
